\section{Introduction}
\label{sec:intro}

Machine unlearning is no longer optional. Privacy regulations
~\cite{gdpr2016,ccpa2018} grant data subjects removal rights, and graph
learning systems--recommendation, fraud detection, citation analysis, and
knowledge-graph retrieval--are squarely in scope. Retraining from scratch
is often prohibitive for production GNNs, so approximate graph unlearning
(GU) methods~\cite{chen2022graph,wu2023gif,cheng2023gnndelete,li2024towards,dong2024idea,zhang2024towards}
trade exact removal for speed via gradient-, influence-, or
partition-based updates.

This trade-off usually assumes benign deletion requests. We show that
the assumption fails once a coalition can choose \emph{which} records to
forget through legitimate APIs. This single degree of freedom exposes a
predictable \textbf{Design-Vulnerability Pattern}: learning-based methods
such as GNNDelete can suffer domino-like F1 collapse (up to $24\%$
absolute drop under degree deletion), influence-based methods are
targeted by gradient-aligned selectors, and partition-based methods show
architectural isolation through Shard Protection. Deletion APIs can
therefore be weaponized unless approximate GU is audited under
worst-case deletion sets.

\paragraph{This paper.}
We conduct the first systematic adversarial audit of approximate GU
under a deletion-API threat model with a calibrated access spectrum
(L0 / L1-structural / L2-\textsc{direct} / L2-\textsc{surrogate};
Table~\ref{tab:access}), treating deletion-set selection as the attack
and post-unlearning F1 as victim utility. To separate family-level trends
from method-specific behavior, we test \textbf{intra-family pairs} drawn
from the upstream OpenGU taxonomy. Our contributions are:

\begin{itemize}[leftmargin=1.2em,topsep=2pt,itemsep=2pt]
  \item \textbf{Threat model and selectors.} We formalize adversarial
        deletion attacks and organize six selectors along a four-tier
        access spectrum keyed on resource cost.
  \item \textbf{Vulnerability Fingerprint.} We project attack-specific F1
        effects onto structural-spread (IM) and gradient-signal (TracIn)
        axes across six GU methods, revealing family-level regimes and
        strong within-family outliers.
  \item \textbf{Alignment and diagnostics.} Mean selected-set degree
        predicts paired F1 effect across $n{=}300$ tuples (Pearson
        $r{=}0.24$, $p{<}10^{-4}$) and within $11/12$
        method$\times$backbone cells; our collateral suite separates the
        $k{=}5$ noise floor, volume effect, attack-specific signal,
        retrain gap, prediction shift, hop decay, and update-detection
        AUC.
  \item \textbf{Defense narrative.} \emph{Shard Protection}---the
        architectural immunity of the partition pair to every selector
        and every budget tested---refines when partition-based GU
        functions as a defense, and the GraphRevoker$\times$GAT
        TracIn breakthrough delineates its boundary.
\end{itemize}

